% Part: computability
% Chapter: recursive-functions
% Section: non-pr-functions

\documentclass[../../../include/open-logic-section]{subfiles}

\begin{document}

\olfileid[tr]{cmp}{rec}{npr}
\olsection{İlkel Özyinelemeli Olmayan Fonksiyonlar}

İlkel özyinelemeli fonksiyonlar, sezgisel olarak hesaplanabilir bütün
fonksiyonları tüketmez. Bütün tek yerli ilkel özyinelemeli fonksiyonları
$f_0,f_1,f_2,\dots$ biçiminde öyle listeleyebildiğimiz sezgisel olarak açık
olmalıdır ki $f_x$ fonksiyonunun $y$ girdisindeki değerini etkin biçimde
hesaplayabilelim. Başka bir deyişle
\[
  g(x,y)=f_x(y)
\]
ile tanımlanan $g(x,y)$ fonksiyonu hesaplanabilirdir. O zaman
\begin{eqnarray*}
  h(x) & = & g(x,x)+1\\
       & = & f_x(x)+1
\end{eqnarray*}
fonksiyonu da hesaplanabilirdir. Her ilkel özyinelemeli $f_i$ fonksiyonu
için $h$ ile $f_i$ değerleri $i$ noktasında farklıdır. Dolayısıyla $h$
hesaplanabilir, fakat ilkel özyinelemeli değildir; $g$ için de aynı şey
söylenebilir. Bu, Cantor'un köşegenleştirme savının ``etkin'' bir sürümüdür.

İlkel özyinelemeli olmayan hesaplanabilir fonksiyonlara daha açık örnekler
verilebilir. $g^n(x)$ gösterimi, toplam $n$ tane $g$ bulunacak biçimde
$g(g(\dots g(x)))$ ifadesini göstersin ve $g_0,g_1,\dots$ fonksiyon dizisi
\begin{eqnarray*}
  g_0(x) & = & x+1\\
  g_{n+1}(x) & = & g_n^x(x)
\end{eqnarray*}
ile tanımlansın. Her $g_n$ fonksiyonunun ilkel özyinelemeli olduğunu
doğrulayabilirsiniz. Birbirini izleyen her fonksiyon öncekinden çok daha hızlı
büyür: $g_1(x)=2x$, $g_2(x)=2^x\cdot x$ ve $g_3(x)$ kabaca $x$ tane $2$'den
oluşan bir üs kulesi gibi büyür. Ackermann--P\'eter fonksiyonu özünde
$G(x)=g_x(x)$ fonksiyonudur ve bunun her ilkel özyinelemeli fonksiyondan daha
hızlı büyüdüğü gösterilebilir.

İlkel özyinelemeli fonksiyonları numaralandırma sorununa dönelim. Her ilkel
özyinelemeli fonksiyona simgesel bir gösterim atadığımızı anımsayın; dolayısıyla
gösterimleri numaralandırmamız yeterlidir. Her $F$ gösterimine bir doğal sayı
$\#(F)$'yi özyinelemeli olarak şöyle atayabiliriz:
\begin{eqnarray*}
  \#(0) & = & \langle0\rangle\\
  \#(S) & = & \langle1\rangle\\
  \#(\Proj{n}{i}) & = & \langle2,n,i\rangle\\
  \#(\fn{Comp}_{k,l}[H,G_0,\dots,G_{k-1}]) & = &
    \langle3,k,l,\#(H),\#(G_0),\dots,\#(G_{k-1})\rangle\\
  \#(\fn{Rec}_l[G,H]) & = & \langle4,l,\#(G),\#(H)\rangle.
\end{eqnarray*}
Burada önceki bölümde verilen kodlar aracılığıyla her sayı dizisinin doğal
bir sayı olarak görülebilmesini kullanıyoruz. Sonuçta her koda bir doğal sayı
atanır. Elbette bazı diziler ve dolayısıyla bazı sayılar hiçbir gösterime
karşılık gelmez. Fakat $i$ böyle bir gösterimi kodluyorsa $f_i$, gösterimi
$i$ ile kodlanan tek yerli ilkel özyinelemeli fonksiyon; aksi durumda sabit
$0$ fonksiyonu olsun. Böylece tek yerli ilkel özyinelemeli fonksiyonları
numaralandırmanın açık bir yolunu elde ederiz.

(Sabit sıfır fonksiyonu gibi bazı fonksiyonlar listede birden fazla kez
görünür. Bu yalnızca kodlamamızın bir yan ürünü değildir; sabit sıfır
fonksiyonunun birden fazla gösterime sahip olmasından da kaynaklanır. Daha
sonra bu tekrarların hesaplanabilir biçimde önlenemeyeceğini göreceğiz.
Örneğin belirli bir gösterimin sabit sıfır fonksiyonunu temsil edip etmediğine
karar veren hesaplanabilir bir fonksiyon yoktur.)

Şimdi $g(x,y)$ fonksiyonunu, az önce açıkladığımız numaralandırmadaki
$f_x(y)$ ile verebiliriz. $g(x,y)$ fonksiyonunun hesaplanabilir olduğunu
nereden biliyoruz? Sezgisel olarak bu açıktır: $g(x,y)$'yi hesaplamak için
önce $x$'in ``paketini açar'', tek yerli bir fonksiyon gösterimi olup
olmadığına bakarız. Öyleyse bu fonksiyonun $y$ girdisindeki değerini
hesaplarız.

\begin{tagblock}{TMs}
\begin{digress}
Bunu yapan bir programın, örneğin Java ya da C++ ile, biraz çalışmayla
yazılabileceğine şimdiden ikna olmuş olabilirsiniz. O zaman sezgisel olarak
hesaplanabilir olan her şeyin bir Turing makinesiyle hesaplanabileceğini
söyleyen Church--Turing tezine başvurabiliriz.

Elbette $g(x,y)$ fonksiyonunun hesaplanabilir olduğunu göstermenin daha
doğrudan yolu, onu hesaplayan bir Turing makinesini açıkça betimlemektir. Bu,
özellikle Church--Turing tezine ve sezgiye başvurmayı önler. Yakında
$g(x,y)$'nin hesaplanabilir olduğunu, bir Turing makinesinde
\emph{benzetilebilen} bir hesaplama modeline, yani özyinelemeli fonksiyonlara
başvurarak gösterecek kadar araç geliştirmiş olacağız.
\end{digress}
\end{tagblock}

\end{document}

